\section{Analysis}\label{sec:analysis}
\begin{proposition}[Exact pooled diagonal moments]\label{prop:moments}
For components with nonnegative masses $a,b$ and population moments $(\bm\mu_a,\bm v_a)$ and $(\bm\mu_b,\bm v_b)$, \cref{eq:poolmean,eq:poolvar} preserves total mass, first moment, and coordinate-wise second moment. Consequently any finite sequence of such merges has the same class mass, mean, and diagonal second moment as the unmerged mixture, up to arithmetic and wire quantization.
\end{proposition}
\begin{proof}
Mass is $m=a+b$ and $m\bm\mu=a\bm\mu_a+b\bm\mu_b$. For each coordinate, substitute \cref{eq:poolvar} into $m(v+\mu^2)$ and expand $(\mu_j-\mu)^2$. Cross terms cancel using $m\mu=a\mu_a+b\mu_b$, leaving $a(v_a+\mu_a^2)+b(v_b+\mu_b^2)$. Induction gives the multi-merge statement.
\end{proof}
Moment preservation is weaker than mixture preservation: distinct distributions share these moments, and a merge may change density valleys or a downstream decision boundary.

\subsection{Packet and Computational Cost}
For dimension $d$, precision $p\in\{2,4\}$ bytes, $C$ represented classes, and $R=\sum_cR_c$ components, the numeric class/component body is
\begin{equation}\label{eq:bytes}
S_{body}=8C+R(8+2pd)\quad\text{bytes}.
\end{equation}
The implementation additionally serializes a four-byte magic value, a four-byte header length, and a JSON header of measured length $H$, so $S=8+H+S_{body}$. Means and variances use precision $p$ on the wire; class and component identifiers are 32-bit integers and masses are float32. Formula \cref{eq:bytes} is analytical and is never presented as measured cumulative communication.

For $n_c$ features, $K_c$ clusters, $I$ Lloyd iterations, and $q$ initializations, center fitting costs $O(qIn_cK_cd)$ without a reservoir, while assignment and exact moment accumulation cost $O(n_cK_cd)$. Storing features costs $O(n_cd)$ in the current implementation; centers and moments cost $O(K_cd)$. A naive greedy merge over $R_c$ components recomputes pair scores and is bounded by $O(dR_c^3)$ for reduction to one, although practical caps are small. Synthesis costs $O(M_cd)$ and edge transformation adds the forward cost of $h$ for the Monte Carlo samples. These are incremental prototype costs; model training and inherited model synchronization remain separate.

\subsection{Conditional Loss Discrepancy}
\begin{assumption}\label{ass:lip}
For fixed parameters $\theta$ and label $c$, feature loss $\ell_\theta(z,c)$ is $L_c$-Lipschitz in the feature space under consideration, and the compared distributions have finite first moments.
\end{assumption}
By Kantorovich--Rubinstein duality,
\begin{equation}\label{eq:lossbound}
|\E_{P_c}\ell_\theta-\E_{Q_c}\ell_\theta|\le L_c\Wone(P_c,Q_c).
\end{equation}
For a class prior $\rho_c$, the total discrepancy is at most $\sum_c\rho_cL_c\Wone(P_c,Q_c)$. This statement is for a fixed model and compatible coordinates. It neither establishes optimization convergence nor says the error is small.

If a boundary map $h$ is $L_h$-Lipschitz, then $W_1(h_\#P,h_\#Q)\le L_hW_1(P,Q)$. In one space, triangle inequalities can separate local fitting, compression, and finite-sample estimation terms. Across $\mathcal Z_{ce}$ and $\mathcal Z_{ec}$ they require the pushforward and $L_h$ factor; directly adding distances from those spaces is invalid. The experiment therefore measures each stage in its own feature space.
